\SourcePage{306}{309}
\section{Atomic energy levels}

In the nonrelativistic approximation, the stationary states of an atom are determined by the Schrödinger equation for a system of electrons moving in the Coulomb field of the nucleus and interacting electrically with one another; electron spin operators do not enter this equation at all. For a system of particles in a centrally symmetric external field, the total orbital angular momentum $L$ and the parity of the state are conserved. Each stationary state of an atom is therefore characterized by a definite value of $L$ and by its parity. Moreover, the coordinate wave functions of stationary states of a system of identical particles have a definite permutation symmetry. We saw in \S~63 that, for a system of electrons, every definite type of permutation symmetry (that is, every Young diagram) corresponds to a definite value of the total spin of the system. Thus every stationary state of an atom is also characterized by the total electron spin $S$.

An energy level with specified $S$ and $L$ is degenerate with respect to the possible directions of the vectors $\mathbf S$ and $\mathbf L$ in space. The degeneracies associated with the directions of $L$ and $S$ are respectively $2L+1$ and $2S+1$. Hence the total degeneracy of a level with specified $L$ and $S$ is $(2L+1)(2S+1)$.

In reality, however, the electromagnetic interaction of electrons includes relativistic effects that depend on their spins. As a result, the energy of an atom depends not only on the magnitudes of $\mathbf L$ and $\mathbf S$, but also on their relative orientation. Strictly speaking, once relativistic interactions are included, neither the orbital angular momentum $L$ nor the spin $S$ of the atom is separately conserved. Only the total angular momentum $\mathbf J=\mathbf L+\mathbf S$ remains conserved. This is a universal exact law following from the isotropy of space for a closed system. The exact energy levels must therefore be characterized by values of the total angular momentum $J$.

\SourcePage{307}{310}
If the relativistic effects are relatively small, as is often the case, they may nevertheless be treated as a perturbation. Under this perturbation, the degenerate level with specified $L$ and $S$ ``splits'' into a series of distinct, closely spaced levels having different values of $J$.

These levels are determined, in first approximation, by the corresponding secular equation (\S~39), while their zeroth-order wave functions are particular linear combinations of the wave functions belonging to the original degenerate level with the given $L$ and $S$.

In this approximation the magnitudes of the orbital angular momentum and spin, though not their directions, may consequently still be regarded as conserved, and the levels may also be characterized by $L$ and $S$.

Thus, through relativistic effects, a level with given $L$ and $S$ splits into levels having different $J$. This splitting is called the \emph{fine structure}, or \emph{multiplet splitting}, of the level. The possible values of $J$ run from $L+S$ to $|L-S|$; a level with given $L$ and $S$ therefore splits into $2S+1$ distinct levels if $L>S$, or $2L+1$ if $L<S$. Each of these remains degenerate with respect to the direction of $\mathbf J$, with degeneracy $2J+1$. It is readily verified that the sum of $2J+1$ over all possible $J$ is, as required, $(2L+1)(2S+1)$.

Atomic energy levels, also called atomic \emph{spectroscopic terms}, are conventionally denoted by symbols analogous to those used for states of individual particles with definite angular momentum (\S~32). States with different total orbital angular momenta $L$ are denoted by capital Latin letters according to
\[
\begin{array}{c|rrrrrrrrrrr}
L&0&1&2&3&4&5&6&7&8&9&10\\
 &S&P&D&F&G&H&I&K&L&M&N
\end{array}
\]
The number $2S+1$, called the \emph{multiplicity} of the term, is placed at the upper left of the symbol (though it should be remembered that this number equals the number of fine-structure components only when $L\geqslant S$).\footnote{For $2S+1=1,2,3,\ldots$, the levels are called singlet, doublet, triplet, and so on.} The value

\SourcePage{308}{311}
of the total angular momentum $J$ is placed at the lower right. Thus ${}^{2}P_{1/2}$ and ${}^{2}P_{3/2}$ denote levels with $L=1$, $S=1/2$, and $J=1/2,3/2$.
